\documentclass[11pt,article,oneside]{memoir}
% \input{vc}

\usepackage{booktabs, amsmath, memoir-article-styles, graphicx, hyperref}

% packages for 'etable' output from R
\usepackage{longtable, makecell, rotating, threeparttable}

\usepackage{pdflscape}

\usepackage[authoryear]{natbib} % citation manager

\usepackage{tikz}
\usetikzlibrary{positioning}

\usepackage[dvipsnames]{xcolor} % Standard suite of named colors

\hypersetup{
    colorlinks=true, % Colored links instead of boxed
    linkcolor={RoyalBlue}, % Color of internal links
    citecolor={RoyalBlue}, % Color of citations
    urlcolor={Cerulean} % Color of URL links
}

\usepackage{float}
\floatplacement{figure}{H} % control figure placement

\setlength{\parindent}{20pt} % Sets paragraph indentation
\setlength{\parskip}{5pt plus 0pt} % Sets paragraph spacing





\title{\bigskip \bigskip Productivity in Manufactured Housing\thanks{I
am grateful to Tyler Pullen for his feedback on an early draft of this
essay.}  }
 

%\author{true}

\author{\Large Colin
Williams\vspace{0.05in} \newline\normalsize\emph{University of
Virginia} \newline\footnotesize \protect\url{chv7bg@virginia.edu}\vspace*{0.2in}\newline }


\date{}


\begin{document}  
\setkeys{Gin}{width=1\textwidth} 	
%\setromanfont[Mapping=tex-text,Numbers=OldStyle]{Minion Pro} 
%\setsansfont[Mapping=tex-text]{Minion Pro} 
%\setmonofont[Mapping=tex-text,Scale=0.8]{Pragmata}
\chapterstyle{article-3} 
\pagestyle{kjh}

\published{March 2026.}

\maketitle



Aggregate data point to a large and decades-long decline in U.S.
construction sector productivity. Value added per worker in 2020 was
roughly 40 percent below its 1970 level, a remarkably poor record for a
major sector \citep{goolsbee_strange_2023}. Conversely, in manufacturing
industries where goods can be produced in factories, productivity has
risen substantially over time \citep{potter_origins_2025}.

What about manufactured housing?

Manufactured homes (MH) are immune to many of the frictions that are
often blamed for poor construction productivity: unpredictable weather,
fragmented local building codes, and the complex coordination of
multiple subcontractors on a construction site. These complications are
absent from the production process for manufactured homes, which are
assembled indoors in a controlled environment and have been regulated by
a single federal code since 1976. However, like site-built construction,
MH remain subject to boom-bust cycles in production and high
transportation costs, limiting the incentive for firms to make large
capital investments that would sit idle during recessions. Thus, the
sector offers a natural test case: are site-specific problems, like
frictions in land assembly or a heavy reliance on contractors, the main
culprits behind poor construction productivity, or is the fundamental
problem with general forces, like demand volatility and transportation
costs, that affect all types of construction?

I start by studying physical output measures. The story is grim:
manufactured housing (MH) shipments per employee fell by roughly half
between the mid-1990s and the mid-2000s and have never fully recovered.
In terms of raw output, MH productivity appears dismal.

Units per worker, however, treats every home equally: a small
single-wide counts the same as a large double-wide. After the demand
collapse of the early 2000s wiped out the low-end chattel-lending
market, the remaining buyers demanded larger, higher-quality units. When
the product mix shifts toward more expensive homes, units per worker can
fall even as dollar-valued productivity holds steady. To better account
for quality, I turn to real value added per employee, from the NBER-CES
database, to measure the dollar value the industry adds per worker,
adjusted for inflation and net of materials costs.\footnote{The NBER-CES
  deflator is based on the Producer Price Index for manufactured
  housing, which tracks prices for specific product categories over
  time. It is not hedonic, so within-product quality improvements ---
  better insulation or appliances at a given price point --- are not
  recognized as output gains. To the extent that such improvements
  occurred, real value added understates true performance.}

This series tells a slightly different story: labor productivity grew
steadily from the late 1950s through the 1990s, then stagnated after
2000 rather than collapsing outright. The divergence between falling
units per worker and flat value added per worker is consistent with a
composition shift toward larger, more expensive homes. At the same time,
however, value added in similar wood product manufacturing industries
continued to grow, suggesting that something specific to the MH industry
caused its relative stagnation.

Measures of total factor productivity (TFP), which strip out the
contribution of capital along with other inputs, suggest an even starker
post-2000 collapse. After years of moderate growth, the NBER-CES TFP
index for the MH sector fell by over 40\% between 1999 and 2015, a stark
divergence from related manufacturing industries.

MH are not immune to the construction productivity puzzle. Indeed, they
appear to have performed even worse than site-built construction,
especially after 2000.

\hypertarget{data}{%
\section{Data}\label{data}}

I combine two data sources to construct a national panel of manufactured
housing productivity from 1958 to 2018.

\textbf{Census Manufactured Housing Survey (MHS).} The MHS provides
annual counts of manufactured home shipments and placements, average
sales prices by unit type (single-wide and double-wide), and state-level
shipment detail. The shipments series extends back to 1959.

\textbf{NBER-CES Manufacturing Industry Database.} The NBER-CES database
provides annual data on employment, value added, shipments, materials
costs, capital stocks, and total factor productivity for all U.S.
manufacturing industries at the six-digit NAICS level from 1958 to 2018
\citep{becker_nber-ces_2021}. Value added is deflated using the
industry-specific shipments price deflator, and the four-factor TFP
index accounts for labor, capital, energy, and materials inputs.

\hypertarget{results}{%
\section{Results}\label{results}}

\hypertarget{physical-output-per-employee}{%
\subsection{Physical output per
employee}\label{physical-output-per-employee}}

I first examine units of housing output per employee for manufactured
housing. MH shipments per worker averaged above five units per employee
through the mid-1990s, then fell sharply to roughly 2.5 units by the
mid-2000s.

\begin{figure}[htbp]
  \centering
  \caption{Physical output per employee}\label{fig:output-pemp}
  \includegraphics[width=\textwidth]{output/output_pemp.pdf}
  \begin{flushleft}
  \begin{footnotesize}
  \emph{  Notes:} MH placements and shipments per employee. Fisher placements are a price-weighted aggregate of single- and double-wide placements.

  \emph{  Source:} Census MHS and NBER-CES Manufacturing Industry Database.
  \end{footnotesize}
  \end{flushleft}
\end{figure}

This decline in physical output is partly the consequence of large and
rapid changes in demand. Financing for MH collapsed in the early 2000s,
eliminating the lowest-income buyers and shifting industry output to
satisfy the remaining customers, with better credit and demand for
larger, higher-quality units. The share of double-wide placements rose
from 50\% of units in 1996 to over 75\% in 2003.

I adjust for this shift by constructing a price-weighted Fisher
aggregate of single- and double-wide placements per employee. This
measure accounts for the shift in the mix of units over time by
weighting the quantity of single- and double-wide placements by their
prices. Importantly, it cannot capture quality improvements within unit
types. The Fisher aggregate shows a similar post-2000 decline, though
the drop is slightly less severe.

\hypertarget{real-value-added-per-employee}{%
\subsection{Real value added per
employee}\label{real-value-added-per-employee}}

I turn next to a measure of real value added per employee from the
NBER-CES database. In contrast to the physical output measure, labor
productivity shows a long upward trend since 1958: from a base of
\$22,000, real value added per worker roughly tripled over the next
forty years to \$63,000 (1997 dollars). Post-2000, value-added stagnated
but did not fall, and by the late 2010s value added per employee
remained near its peak.

\begin{figure}[htbp]
  \centering
  \caption{Real value added per employee}\label{fig:va-pemp}
  \includegraphics[width=\textwidth]{output/va_pemp.pdf}
  \begin{flushleft}
  \begin{footnotesize}
  \emph{  Notes:} Value added deflated by shipments price index (1997 =
  1.0). Comparison series are aggregates for other NAICS 321 (wood product manufacturing) industries.
  
  \emph{  Source:} NBER-CES Manufacturing Industry Database.
  \end{footnotesize}
  \end{flushleft}
\end{figure}

At the same time, however, the industry performed much worse than
comparable manufacturing sectors. Wood product manufacturing industries
grew steadily through the 2000s, ending the 2010s with nearly double the
real value added per employee of MH despite starting at a similar level
in the
1990s.\footnote{Falling input prices, such as lumber, could in principle affect measured real value added per worker. However, other wood product manufacturing industries faced the same input price changes, so this force alone cannot explain the MH-specific divergence.}

The most likely explanation for the divergence between raw output and
value-added measures is that the composition of output shifted towards
higher-quality, more expensive units. This shift is consistent with
gradual increases to minimum standards in the federal HUD code,
especially for wind and energy standards during the mid-1990s, which
raised the quality of the product \citep{treaster_considered_2017}.

\hypertarget{total-factor-productivity}{%
\subsection{Total factor productivity}\label{total-factor-productivity}}

Finally, I plot the NBER-CES four-factor TFP index, which adjusts for
labor, capital, energy, and materials inputs, for manufactured housing
and Domar-weighted aggregates of other manufacturing
industries.\footnote{Aggregate TFP series for comparison groups are
  constructed following \citet{domar_measurement_1961}. For each group,
  the annual TFP growth rate is a weighted average of constituent
  industries' four-factor TFP growth rates, where the weight on each
  industry is its lagged gross output (shipments) as a share of the
  group's lagged aggregate value added. These Domar weights sum to
  greater than one, reflecting the fact that an industry's productivity
  gain raises the effective output of downstream users of its products.
  The weighted growth rates are cumulated from a base of 1.0 in 1997.}
The story is quite similar to value added: MH TFP rose steadily between
1960 and 1999, then declined precipitously to 0.54 by 2013. The partial
recovery to 0.65 by 2016 still leaves TFP well below its level in 1960.

\begin{figure}[htbp]
  \centering
  \caption{Total factor productivity}\label{fig:tfp}
  \includegraphics[width=\textwidth]{output/tfp.pdf}
  \begin{flushleft}
  \begin{footnotesize}
  \emph{  Notes:} Four-factor TFP index using labor, capital, energy, and
  materials inputs. Base year 1997 = 1.0. Comparison series are
  Domar-weighted aggregates of annual industry TFP growth for
  other NAICS 321 (wood product manufacturing) industries.

  \emph{  Source:} NBER-CES Manufacturing Industry Database.
  \end{footnotesize}
  \end{flushleft}
\end{figure}

By contrast, other wood product manufacturing industries show no break
in trends around 2000. MH's productivity collapse was not part of a
broad manufacturing phenomenon.

The juxtaposition of Figures \ref{fig:va-pemp} and \ref{fig:tfp} implies
that labor productivity recovered only through capital deepening: firms
invested in more capital per worker, but the efficiency with which they
combine all inputs deteriorated dramatically. Manufactured housing
factories, despite their controlled environment, standardized processes,
and federal regulation, have not escaped the forces holding down
construction productivity.

\hypertarget{discussion}{%
\section{Discussion}\label{discussion}}

The TFP collapse after 1998 coincides almost exactly with the demand
shock that hit manufactured housing when the subprime chattel-lending
channel dried up. MH shipments fell from roughly 375,000 units per year
at their late-1990s peak to fewer than 80,000 by the mid-2000s, a
massive decline from which the industry has never recovered. One
interpretation of the TFP decline is that it reflects lost economies of
scale rather than technological regress.

The scale interpretation has both within-plant and market-level
dimensions. Within each factory, high fixed costs --- the plant itself,
jigs, tooling, and quality-control infrastructure --- must be spread
over output. When volumes collapsed, factories that once operated at
efficient scale were left with excess capacity, mechanically depressing
measured productivity. But the geographic fragmentation of the industry
compounds this problem. Transport costs for finished manufactured homes
are steep (about \$13--14 per mile), so each factory serves a limited
radius \citep{jensen_manufactured_2024}. Unlike a typical manufacturing
industry, MH producers cannot respond to falling demand by consolidating
into fewer, larger plants: a factory in Alabama cannot absorb the
customers of a shuttered plant in Oregon. Some regional markets may
simply be unable to sustain a factory at efficient scale, with no
prospect of consolidation.

The severity of the demand shock also makes it unlikely that standard
reallocation forces offset the productivity decline. In a Melitz-style
framework, negative demand shocks should induce the least productive
firms to exit, raising average productivity among survivors. The fact
that MH TFP fell sharply despite substantial industry contraction
suggests that the loss of scale economies dominated any cleansing effect
from the exit of low-productivity plants.

A lack of scale economies may also explain the poor performance of
site-built construction. \citet{damico_why_2024} argue that fragmented
local regulations prevent large, efficient builders from producing
standardized homes on large greenfield sites, a mode of production that
was historically more common and which contributed to the industry's
productivity growth in the years after World War II.

At the same time, the scale interpretation has a silver lining. If the
post-2000 productivity collapse reflects lost volume rather than
technological regress, then the production knowledge still exists and
factories \emph{can} be efficient --- they just need demand. Reforms to
titling and financing --- for example, making manufactured homes
eligible for conventional mortgage products --- could expand the buyer
pool and set off a virtuous cycle: higher volumes would lower unit
costs, which would lower prices, which would attract additional buyers
at the margin. Because the industry is likely operating on the steep
portion of its average cost curve, even modest demand increases could
yield disproportionate productivity gains.

The geographic structure of the industry, however, limits how far this
logic can run. Because high transport costs confine each factory to a
limited radius, national demand growth must translate into
\emph{regionally distributed} demand growth to restore scale at
individual plants. Full recovery to 1990s volumes appears unlikely given
secular headwinds --- falling fertility, regulatory fragmentation, and
the permanent loss of the subprime chattel-lending channel. For both
site-built and manufactured homes, structural features of the industry
limit the size of the market, prevent learning-by-doing, and raise the
cost of housing.

\bibliographystyle{chicago}
\bibliography{manufactured-productivity.bib}

% original bibliography code
% % %%
\end{document}